% -*- compile-command: "pdflatex -shell-escape Hendrix-comprog-reference.tex" -*-

\documentclass[10pt]{book}

\usepackage[margin=1in]{geometry}
\usepackage{fancyhdr}
\usepackage{amsmath}

\pagestyle{fancy}

% \lhead{\leftmark}
% \rhead{\rightmark}

% \lfoot{Hendrix Programming Team Reference}
% \rfoot{\thepage}

% TODO: make fancy header with page number and section etc.

\usepackage[T1]{fontenc}
\usepackage[utf8]{inputenc}
\usepackage{hyperref}
\usepackage{xcolor}

\hypersetup{
    colorlinks,
    linkcolor={green!50!black},
    citecolor={blue!50!black},
    urlcolor={blue!80!black}
}

\usepackage{titlepic}
\usepackage{graphicx}
\graphicspath{{images/}}

\usepackage[newfloat]{minted}
\usepackage{etoolbox}
\usepackage{mdframed}
\usepackage{booktabs}
\usepackage{xspace}
\usepackage{pifont}
\usepackage{xparse}

\usepackage{algorithm, algorithmicx}
\usepackage[noend]{algpseudocode}

\usepackage{caption}
\DeclareCaptionLabelFormat{algnonumber}{Algorithm}
\captionsetup[algorithm]{labelformat=algnonumber}

%%%%%%%%%%%%%%%%%%%%%%%%%%%%%%%%%%%%%%%%%%%%%%%%%%
%% Inline code

\NewDocumentCommand{\code}{O{}O{}m}{\inputminted[linenos=true,mathescape,firstline=#1,lastline=#2,autogobble]{java}{code/#3}}
\newcommand{\pycode}[1]{\inputminted[linenos=true,mathescape]{python}{code/#1}}

%%%%%%%%%%%%%%%%%%%%%%%%%%%%%%%%%%%%%%%%%%%%%%%%%%
%% Links

\newcommand*{\fulllink}[1]{\hyperref[{#1}]{\nameref*{#1}~(\S\ref*{#1}, page~\pageref*{#1})}}
\newcommand*{\link}[1]{\hyperref[{#1}]{(\S\ref*{#1}, page~\pageref*{#1})}}

\newcommand{\javalogo}{\includegraphics[height=0.9\baselineskip]{Java}}
\newcommand*{\javadoclink}[2]{\href{https://docs.oracle.com/javase/10/docs/api/java/#1/#2.html}{\texttt{#2}}}
\newcommand*{\javadoc}[2]{\javalogo\ \javadoclink{#1}{#2}}

%%%%%%%%%%%%%%%%%%%%%%%%%%%%%%%%%%%%%%%%%%%%%%%%%%
%% Kattis

\newcommand{\kattislogo}{\raisebox{-0.2em}{\includegraphics[height=0.9\baselineskip]{Kattis}}}

\newcommand{\kattislist}[1]{%
  \def\nextitem{\def\nextitem{, }}
  \renewcommand*{\do}[1]{\nextitem\kattislink{##1}}
  \docsvlist{#1}
}

\newcommand{\kattis}[1]{%
  \begin{mdframed}
    \kattislogo
    \kattislist{#1}
  \end{mdframed}
}

\newcommand{\kattislink}[1]{\href{https://open.kattis.com/problems/#1}{\texttt{#1}}}

\newcommand{\inlinekattis}[1]{\kattislogo\!\!\kattislist{#1}\!}

%%%%%%%%%%%%%%%%%%%%%%%%%%%%%%%%%%%%%%%%%%%%%%%%%%
%% Todos

\newif\iftodos
\todostrue
\newcommand{\todo}[1]{\iftodos\textcolor{red}{[TODO: #1]}\fi}

%%%%%%%%%%%%%%%%%%%%%%%%%%%%%%%%%%%%%%%%%%%%%%%%%%
%% Warning

\newenvironment{warning}
{\par\begin{mdframed}[linewidth=2pt,linecolor=red]%
    \begin{list}{}{\leftmargin=1cm
        \labelwidth=\leftmargin}\item[\Large\ding{43}]}
    {\end{list}\end{mdframed}\par}

%%%%%%%%%%%%%%%%%%%%%%%%%%%%%%%%%%%%%%%%%%%%%%%%%%
%% Misc

\newcommand{\eg}{\emph{e.g.}\xspace}
\newcommand{\ie}{\emph{i.e.}\xspace}

%%%%%%%%%%%%%%%%%%%%%%%%%%%%%%%%%%%%%%%%%%%%%%%%%%%%%%%%%%%%%%%%%%%%%%

\begin{document}

\title{Hendrix Programming Team Reference}
\titlepic{\includegraphics[width=3in]{Hendrix-logo}}
\maketitle

\tableofcontents
\newpage

\chapter{Limits}

As a rule of thumb, you should assume about $10^8$ (= 100 million)
operations per second.  If you can think of a straightforward brute
force solution to a problem, you should check whether it is likely to
fit within the time limit; if so, go for it!  Some problems are
explicitly written to see if you will recognize this.  If a brute
force solution won't fit, the input size can help guide you to search
for the right algorithm running time.

Example: suppose a problem requires you to find the length of a
shortest path in a weighted graph.
\begin{itemize}
\item If the graph has $|V| = 400$ vertices, you should use
  Floyd-Warshall \link{sec:floydwarshall}: it is the easiest to code and takes $O(V^3)$ time
  which should be good enough.
\item If the graph has $|V| = 4000$ vertices, especially if it doesn't
  have all possible edges, you can use Dijkstra's algorithm
  \link{sec:dijkstra}, which is $O(E \log V)$.
\item If the graph has $|V| = 10^5$ vertices, you should look for some
  special property of the graph which allows you to solve the problem
  in $O(V)$ or $O(V \log V)$ time---for example, perhaps the graph is
  a tree \link{sec:graph-basics}, so you can run a BFS/DFS \link{sec:dfs}
  to find the unique path and then add up the weights.  An input size of
  $10^5$ is a common sign that you are expected to use an $O(n \lg n)$
  or $O(n)$ algorithm---it's big enough to make $O(n^2)$ too slow but
  not so big that the time to do I/O makes a big difference.
\end{itemize}

\begin{table*}[h]\centering
  \renewcommand{\arraystretch}{1.3}
  \begin{tabular}{@{}lll@{}}\toprule
    $n$ & Worst viable running time & Example \\
    \midrule
    $11$ & $O(n!)$ & Generating all permutations \link{sec:combinatorics} \\
    $25$ & $O(2^n)$ & Generating all subsets \link{chap:bittricks} \\
    $100$ & $O(n^4)$ & Some brute force algorithms \\
    $400$ & $O(n^3)$ & Floyd-Warshall \link{sec:floydwarshall} \\
    $10^4$ & $O(n^2)$ & Testing all pairs \\
    $10^6$ & $O(n \lg n)$ & BFS/DFS; sort+greedy \\
    \bottomrule
  \end{tabular}
\end{table*}

\kattis{bing, transportationplanning, dancerecital, prozor,
  rectanglesurrounding, weakvertices}

\begin{itemize}
\item $2^{10} = 1024 \approx 10^3$.
\item One \texttt{int} is 32 bits = 4 bytes. So \emph{e.g.} an array
  of $10^6$ \texttt{int}s requires $< 4$ MB---no big deal since the
  typical memory limit is $1$ GB.  Don't be afraid to make arrays with
  millions of elements!
\item \texttt{int} holds 32 bits; the largest \texttt{int} value is
  \verb|Integer.MAX_VALUE| $ = 2^{31} - 1$, a bit more than
  $2 \cdot 10^9$.
\item \texttt{long} holds 64 bits; the largest
  \texttt{long} value is \verb|Long.MAX_VALUE| $ = 2^{63} - 1$, a bit
  more than $9 \cdot 10^{18}$. To write literal long values you can
  add an \texttt{L} suffix, as in \texttt{long x = 1234567890123L;}.
\item If you need larger values, use \fulllink{sec:bigint} or just use
  Python~\link{chap:python}; see also \fulllink{sec:combinatorics}.
\end{itemize}

\kattis{different}

\chapter{Java Reference}

\section{Template}

\code{java/Template.java}

\section{Scanner}

\javadoc{util}{Scanner} is relatively slow but should usually be sufficient
for most purposes.  If the input or output is relatively large (> 1MB)
and you suspect the time taken to read or write it may be a hindrance,
you can use \fulllink{sec:fastio}.

\begin{warning}
  Be sure to read the warning in the comment below about calling
  \texttt{nextLine()} after \texttt{nextInt()} and the like!
\end{warning}

\code{java/ScannerExample.java}

\section{Math}

The standard \javadoc{lang}{Math} class contains useful standard
mathematical constants and operations.  All are \texttt{static}, so
they can be accessed by prefixing their names with \texttt{Math.}, \ie
\texttt{Math.cos}.
\begin{itemize}
\item Constants \texttt{E} and \texttt{PI} represent (floating-point
  approximations of) $e$ and $\pi$.
\item \texttt{abs} finds the absolute value.
\item \texttt{min} and \texttt{max} find the min or max of two
  values. A common trick for saving a bit of typing is to use
  something like
  \begin{minted}{java}
    m = Math.max(m, val);
  \end{minted}
  if you need \texttt{m} to accumulate the maximum of a set of values.
\item \texttt{round} rounds a floating-point number to the nearest
  integer. \texttt{ceil} and \texttt{floor} round up and down,
  respectively.  Note that whereas \texttt{round}
  returns a \texttt{long} when given a \texttt{double}, for some
  reason \texttt{ceil} and \texttt{floor} both return
  \texttt{double}, so you may need to cast the results:
  \begin{minted}{java}
    double x = ...
    long n = (long)Math.floor(x);
  \end{minted}
\item \texttt{exp(x)} computes $e^x$.  \texttt{log(x)} computes $\ln x$.
\item \texttt{sqrt} computes the square root.  \texttt{hypot(x,y)}
  computes $\sqrt{x^2 + y^2}$.
\item \texttt{pow(a, b)} computes $a^b$.
\item \texttt{sin}, \texttt{cos}, \texttt{tan}, \texttt{acos},
  \texttt{asin}, \texttt{atan} do what you would expect. Note also
  \texttt{atan2(y,x)} which returns an angle $\theta$ such that
  it converts rectangular coordinates $(x,y)$ into polar coordinates
  $(r, \theta)$.  This is almost like \texttt{atan(y/x)} except that
  it avoids division by zero and handles all four quadrants properly.
\item \texttt{toDegrees} and \texttt{toRadians} convert angles from
  radians to degrees and degrees to radians, respectively.
\end{itemize}


\section{String}

\kattis{battlesimulation, bing, connectthedots, itsasecret, shiritori,
suffixarrayreconstruction}

The \javadoc{lang}{String} type can be used in Java to represent
sequences of characters.  Some useful \texttt{String} methods include:
\begin{itemize}
\item concatenation (\texttt{+})
\item \texttt{substring(i)} yields the substring starting at index
  \texttt{i} up to the end of the string
\item \texttt{substring(i,j)} yields the substring starting at
  \texttt{i} (inclusive) and ending \emph{just before} \texttt{j}
  (same as Python slices).
\item \texttt{charAt(i)} yields the \texttt{char} at index \texttt{i}.
\item \texttt{toCharArray()} converts to a \texttt{char[]}, which can
  be convenient if you need to do a lot of indexing (\texttt{[i]}
  instead of \texttt{charAt(i)})
\item \texttt{split(String)} splits a string into a \texttt{String[]}
  of pieces between occurrences of the splitting string.
\item \texttt{endsWith(String)}, \texttt{startsWith(String)},
  \texttt{indexOf(String)}, and \texttt{replace(...)} can occasionally be useful.
\end{itemize}

Below is shown a solution to \inlinekattis{sumoftheothers}, which uses
\texttt{split} followed by \texttt{Integer.parseInt} to read the
integers on each line (necessary in this case because the input does
not specify how many integers will be on each line, although this is atypical).

\code{java/sumOfTheOthers.java}

\texttt{String}s are immutable, which means in particular that
concatenation has to allocate an entirely new \texttt{String} and copy
both arguments.  Hence repeatedly appending individual characters to
the end of a \texttt{String} takes $O(n^2)$ time, since the entire
string must be copied with each append operation.  In this situation,
either pre-allocate a sufficiently large \texttt{char[]}, or use
the \javadoc{lang}{StringBuilder} class.

\section{StringBuilder}

\kattis{itsasecret, joinstrings}

\javadoc{util}{StringBuilder} is a \emph{mutable} string class which
supports efficient append and modification operations.  If you need to
build up a long string by incrementally appending text bit by bit, you
should use \texttt{StringBuilder} instead of using \texttt{String}
directly.  \texttt{StringBuilder} also supports a \texttt{reverse()}
method (unlike \texttt{String}).

As a simple example, the below code prints \texttt{0291817161514131211101987654321}.

\begin{minted}{java}
  StringBuilder sb = new StringBuilder();
  for (int i = 1; i <= 20; i++) {
      sb.append("" + i);
  }
  System.out.println(sb.reverse());
\end{minted}

\section{Arrays}

\kattis{falcondive,freefood,traveltheskies}

The basic syntax for creating a primitive array in Java is, for example,
\begin{minted}{java}
int[] array = new int[500];
\end{minted}

Some tips and tricks:
\begin{itemize}
\item Array indexing starts at 0; however, problems sometimes index
  things from $1 \dots n$.  In such a situation it is usually a good
  idea to simply create an array with one extra slot and leave index 0
  unused.  The alternative (fiddling with indices by subtracting and
  adding 1 in the right places) is quite error-prone.
\item You can initialize an entire array to a given value using
  \texttt{Arrays.fill(array, value)}.
\item If you only want to initialize part of an array, use
  \texttt{Arrays.fill(array, fromIndex, toIndex, value)} to fill the
  array from \texttt{fromIndex} (inclusive) up to \texttt{toIndex}
  (exclusive).
\item You can sort the contents of an array in-place using
  \href{https://docs.oracle.com/javase/10/docs/api/java/util/Arrays.html#sort(int\%5B\%5D)}{\texttt{Arrays.sort}};
  see \fulllink{sec:sorting}.
\item You can use \texttt{Arrays.binarySearch(array, key)} to look for
  \texttt{key} within a sorted \texttt{array}.  Read
  \href{https://docs.oracle.com/javase/10/docs/api/java/util/Arrays.html#binarySearch(int\%5B\%5D,int)}{the
    documentation} to make sure you understand how to interpret the
  return value.  See also \fulllink{sec:binary-search}.
\item Other methods from the \javadoc{util}{Arrays} class may also
  occasionally be useful.
\item To iterate over the items in an array, you can use foreach
  syntax:
  \begin{minted}{java}
    for (int item : array) {
      // do something with i
    }
  \end{minted}
\end{itemize}

\section{ArrayList}

\javadoc{util}{ArrayList} represents a standard dynamically-extensible
array, doubling the underlying storage when it runs out of space so
that appending takes $O(1)$ amortized time.  The \texttt{add},
\texttt{get}, \texttt{set}, \texttt{size}, and \texttt{isEmpty}
methods are useful, in addition to the ability to iterate over the
elements in order.  Avoid methods such as \texttt{contains},
\texttt{indexOf}, \texttt{remove}, and the version of \texttt{add}
that takes an index, all of which take linear time.  (If you think you
want any of these methods, it's probably a sign that you ought to be
using a different data structure.)

If you need to store a list/array and you know in advance exactly how
much storage space you will need, then prefer using a primitive array
which has less overhead as well as more concise syntax.  On the other
hand, if you want to be able to dynamically extend a list by appending
new elements to the end, use \texttt{ArrayList}.  (If you want to be
able to dynamically extend a list on \emph{both} ends, use an
\texttt{ArrayDeque} \link{sec:queue}.)

\begin{minted}{java}
ArrayList<Integer> lst = new ArrayList<>();
lst.add(3); lst.add(19); lst.add(6);
System.out.println(lst.get(2));   // prints 6
lst.set(1, 12);                   // changes 19 to 12
int sum = 0;
for (Integer i: lst) {            // iterate through all items
    sum += i;
}
System.out.println(sum);          // prints 3 + 12 + 6 = 21
\end{minted}

\section{Stack}

\kattis{backspace, delimitersoup, islands, pairingsocks, reservoir, restaurant, symmetricorder,
  throwns, zagrade}

\javadoc{util}{Stack} provides a generic stack implementation with
$O(1)$ operations.  Standard methods include \texttt{isEmpty},
\texttt{push}, \texttt{pop}, \texttt{peek}, and \texttt{size}.  The
code below shows a sample solution to \inlinekattis{backspace} using
\javadoc{util}{Stack} (and \javadoc{util}{StringBuilder}).

\code{java/backspace.java}

Stacks are often used in implementing DFS \link{sec:dfs} as well as
dealing with parentheses, or nesting more generally
(\inlinekattis{pairingsocks}, \inlinekattis{islands},
\inlinekattis{reservoir}).

\section{Queue/ArrayDeque}
\label{sec:queue}

\kattis{eenymeeny, brexit, coconut, ferryloading4, integerlists, shuffling}

\javadoc{util}{Queue}, unlike \javadoc{util}{Stack}, is not a class
but an interface.  There are several classes implementing the
\texttt{Queue} interface, but the best in the context of competitive
programming is probably \javadoc{util}{ArrayDeque}, which in fact
implements a \emph{double-ended queue} or \emph{deque}, providing
$O(1)$ amortized addition and removal from both ends.

The \texttt{add} and \texttt{remove} methods implement enqueueing and
dequeueing.  To access both ends, use \texttt{addFirst},
\texttt{addLast}, \texttt{removeFirst}, and \texttt{removeLast}, all
of which run in $O(1)$ amortized time. (\texttt{add} is the same as
\texttt{addLast} and \texttt{remove} is the same as
\texttt{removeFirst}.)

Queues are very commonly used in implementing \fulllink{sec:bfs} and
in simulations of various sorts (for examples of the latter, see the
selection of problems above).

\todo{Sample code using \texttt{Queue}/\texttt{ArrayDeque}}

\section{Comparator}
\label{sec:comparator}

A \javadoc{util}{Comparator} is used to specify a custom ordering on
some type, potentially different than its ``natural'' ordering.
Typically a \texttt{Comparator} can be passed as an optional argument
to things that require an ordering.  For example, given an
\texttt{ArrayList<Integer> arr}, one can use
\texttt{Collections.sort(arr)} to sort it in increasing numeric order.
If a different order is wanted, one can pass a \texttt{Comparator} as
the second argument to \texttt{sort}, as in
\begin{minted}{java}
Collections.sort(arr, Collections.reverseOrder());
\end{minted}
to sort in descending order, or
\begin{minted}{java}
Collections.sort(arr, (i,j) -> q[i] - q[j]);
\end{minted}
to sort a list of \emph{indices} by the corresponding value in an
array \texttt{q}.  A \texttt{Comparator} can also be used as an extra
argument to the constructor when creating a data structure that
depends on ordering, such as a \texttt{PriorityQueue},
\texttt{TreeSet}, or \texttt{TreeMap}.

\todo{Constructing Comparators via lambda; constructing via things
  like \texttt{comparing}, \texttt{thenComparing}.
  \texttt{Collections.reverseOrder()}.}

\section{PriorityQueue}
\label{sec:pq}

\kattis{bank, ferryloading3, guessthedatastructure, knigsoftheforest, vegetables}

A \javadoc{util}{PriorityQueue} allows adding new elements
(\texttt{add}) and removing the \textbf{minimum} element
(\texttt{remove}), both in $O(\lg n)$ time.  \texttt{peek} can also
be used to get the minimum in $O(1)$ without removing it.  Priority
queues are commonly used in Dijkstra's algorithm \link{sec:dijkstra},
event-based simulations (\inlinekattis{ferryloading3}), and generally
any situation where we need to do an ``online sort'', that is, we need
to get items in order from smallest to biggest, but more items may
continue to arrive/be generated as we go.

Methods you should \emph{not} use with \texttt{PriorityQueue} include
\texttt{remove(Object)} and \texttt{contains(Object)}, which take
linear time.

The default constructor makes an empty min-PQ.  If you want to use a
different ordering, there is another constructor which takes a
\javadoc{util}{Comparator}.
\begin{itemize}
\item For example, if you want a \textbf{max} priority queue, where
  \texttt{remove()} yields the largest element, write something like
  \begin{minted}{java}
    PriorityQueue<Integer> pq = new PriorityQueue<>(Collections.reverseOrder());
  \end{minted}
\item If you want some other ordering, you can also use a lambda to
  construct a \texttt{Comparator} on the fly, for example:
  \begin{minted}{java}
    PriorityQueue<Integer> pq = new PriorityQueue<>((a,b) -> dist[a] - dist[b]);
  \end{minted}
\end{itemize}

Traditional presentations of priority queues often have a
\emph{decrease key} operation which can decrease the priority of an
item (or an \emph{adjust key} operation which can arbitrarily change
the priority) and reestablish the data structure invariants in
$O(\lg n)$ time; this operation is used, for example, in implementing
Dijkstra's algorithm efficiently \link{sec:dijkstra}.  However, the
Java \texttt{PriorityQueue} class has no such method.  One workaround
is to simply call \texttt{remove} and then \texttt{add} so the item
gets re-added with the new priority.  However, \texttt{remove} takes
linear time, so this is not ideal, although in many cases it is still
good enough.  For those (relatively rare) cases when an $O(\lg n)$
decrease key operation is truly essential, see \fulllink{sec:adj-pq}.

\section{Set}

\kattis{boatparts, bookingaroom, engineeringenglish,
  whatdoesthefoxsay, securedoors, bard, control}

The \javadoc{util}{Set} interface represents a collection of items
where each item occurs at most once. Operations supported by all
\texttt{Set}s include \texttt{add}, \texttt{remove}, \texttt{contains},
\texttt{size}/\texttt{isEmpty}.

There are two main classes implementing \texttt{Set}:

\begin{itemize}
\item \javadoc{util}{HashSet} is implemented using a dynamically
  expanding hash table.  It features $O(1)$ \texttt{add},
  \texttt{remove}, and \texttt{contains}.
\item \javadoc{util}{TreeSet} is implemented using a balanced binary
  tree (a red-black tree, in fact), and supports \texttt{add},
  \texttt{remove}, and \texttt{contains} in guaranteed $O(\lg n)$
  time.  However, it has several other advantages over a \texttt{HashSet}:
  \begin{itemize}
  \item Since the elements are stored in order in the tree, iterating
    over a \texttt{TreeSet} is guaranteed to yield the items in order
    from smallest to biggest, whereas iterating over a
    \texttt{HashSet} yields the items in an arbitrary order.  For
    example, if you want to remove duplicates from a set of items and
    then print them out in order, you might as well just throw them
    all into a \texttt{TreeSet} instead of putting them in a
    \texttt{HashSet} and then sorting (\inlinekattis{crowdcontrol}).
    (And either one is probably going to be faster than sorting and
    \emph{then} removing duplicates.)
  \item If you need to put objects of a custom class into a set, it is
    sometimes easier to implement \texttt{Comparable} for your class
    and use a \texttt{TreeSet} than it is to override
    \texttt{hashCode} and use a \texttt{HashSet}.  The $O(\lg n)$
    difference is rarely, if ever, going to be the difference between
    AC and TLE, so you should use whichever approach will be easier to
    code.
  \item \texttt{TreeSet} also supports the \javadoc{util}{OrderedSet}
    and \javadoc{util}{NavigableSet} interfaces, which provide
    additional methods like \texttt{first} and \texttt{last} (return
    the smallest or largest element in the set), \texttt{headSet} and
    \texttt{tailSet} (return the subset of all items less or greater
    than a specified element), and \texttt{floor}, \texttt{ceiling},
    \texttt{lower}, and \texttt{higher} (find the first item in the
    set less/greater than a specified value).  This last set of
    methods can be especially useful for some types of problems.
    \kattis{closestsums, platforme, baloni, excellentengineers}
  \end{itemize}

\end{itemize}

There is also a \javadoc{util}{LinkedHashSet} class which in addition
to providing all the same features as a \texttt{HashSet}, also
remembers the order in which the items were added; iterating over the
set is guaranteed to yield the items in this order.  This is a bit
more sophisticated than simply keeping an \texttt{ArrayList} and a
\texttt{HashSet} side-by-side, in particular because a
\texttt{LinkedHashSet} still supports $O(1)$ removal.  We currently do
not know of any example problems which can be solved most easily using
a \texttt{LinkedHashSet}, but it never hurts to be prepared!

\section{Map}
\label{sec:map}

\kattis{awkwardparty, administrativeproblems, snowflakes, pizzahawaii, snowflakes}

The \javadoc{util}{Map} interface represents a dictionary data
structure associating keys to values.  Supported operations include
\texttt{get(K)}, \texttt{put(K,V)}, \texttt{containsKey(K)},
\texttt{remove(K)}, and \texttt{size/isEmpty}.

As with \texttt{Set}, there are two main classes implementing
\texttt{Map}:

\begin{itemize}
\item \javadoc{util}{HashMap} is implemented using a hash table, and
  allows $O(1)$ \texttt{get}, \texttt{put}, \texttt{remove}, and
  \texttt{containsKey}.

\item \javadoc{util}{TreeMap} is implemented using a balanced binary
  tree.  Many of the same comments apply as for \texttt{TreeSet}:
  \begin{itemize}
  \item All operations run in worst-case $O(\lg n)$ time.
  \item Iterating over the keys in the map is guaranteed to return
    them in order from smallest to biggest.
  \item \texttt{TreeMap} also implements the \javadoc{util}{SortedMap}
    interface (allowing one \emph{e.g.} to access the first or last
    key or to get a submap of all the key/value pairs which lie in
    between certain keys) and \javadoc{util}{NavigableMap} interface
    (which lets you find the closest keys and values which are
    smaller/bigger than a given query key).
  \end{itemize}
\end{itemize}

For the purposes of programming contests, \texttt{TreeMap} and
\texttt{HashMap} are basically interchangeable.  \texttt{HashMap} is
faster in theory but a factor of $\lg n$ is not that much, and
\texttt{HashMap} has its own overhead costs.

\begin{itemize}
\item To iterate over the keys of a map, use \texttt{keySet}:
  \begin{minted}{java}
    for (K key : map.keySet()) {
      ...
    }
  \end{minted}
\item To iterate over the values, use \texttt{values}:
  \begin{minted}{java}
    for (V val : map.values()) {
      ...
    }
  \end{minted}
\item You can also iterate over both at once:
  \begin{minted}{java}
    for (Map.Entry<K,V> e : map.entrySet()) {
      ... e.getKey() ... e.getValue() ...
    }
  \end{minted}
\end{itemize}

\section{BigInteger} \label{sec:bigint}

\kattis{basicremains}

\javadoc{math}{BigInteger} represents arbitrarily large integer
values, though it's not actually needed all that often You can import
it as \texttt{java.math.BigInteger}.

\begin{itemize}
\item To turn an \texttt{int} or \texttt{long} into a
  \texttt{BigInteger}, use \texttt{BigInteger.valueOf}.
\item \texttt{BigInteger} also provides the constants
  \texttt{BigInteger.ZERO}, \texttt{.ONE}, \texttt{.TWO}, and
  \texttt{.TEN}.
\item To do arithmetic, use methods such as \texttt{add},
  \texttt{subtract}, and \texttt{multiply}, which all return their
  result as a new \texttt{BigInteger}:
  \begin{minted}{java}
    BigInteger a = ...
    BigInteger b = ...
    BigInteger c = a.add(b.multiply(BigInteger.TEN));  // c = a + b*10
  \end{minted}
\item \texttt{BigInteger} can occasionally be useful in scenarios
  other than representing large numbers, because of its useful utility
  methods such as \texttt{gcd} and its ability to convert between
  different bases.
\end{itemize}

\section{Sorting}
\label{sec:sorting}

\begin{itemize}
\item You can sort an array using \texttt{Arrays.sort(array)}.
\item You can sort a list-like collection (such as an
  \texttt{ArrayList}) using \texttt{Collections.sort(list)}.
\item These \texttt{sort} methods also take an optional
  \fulllink{sec:comparator}.
\end{itemize}

You should rarely, if ever, have to code your own sorting algorithms;
the exception is when you need to \emph{enhance} a sorting algorithm
to keep track of some extra information while sorting (for example,
look up the divide-and-conquer algorithm for counting inversions).

It is relatively common that one needs to sort according to some key,
but carry along additional information with the keys.  The easiest way
to do this is to make a small class to contain all the relevant
information, then implement the \texttt{Comparable} interface.
\begin{minted}{java}
class Person implements Comparable<Person> {
    int age; String name;
    public Person(int age, String name) { this.age = age; this.name = name; }
    int compareTo(Person p) { return this.age - p.age; }
}

ArrayList<Person> people = ...
Collections.sort(people);  // sort by age, carrying names along too
\end{minted}

\section{Fast I/O} \label{sec:fastio}

\kattis{avoidland, cd, ozljeda}

Typically ACM ICPC problems are designed so \texttt{Scanner} and
\texttt{System.out.println} are fast enough to read and write the
required input and output within the time limits.  However, these are
relatively slow since they are unbuffered (every single read and write
happens immediately).  Occasionally it can be useful to have faster
I/O; indeed, a few problems on Kattis cannot be solved in Java without
using this.

\begin{warning}
  Be sure to call \texttt{flush()} at the end of your program
  or else some output might be lost!
\end{warning}

\code{java/Kattio.java}

% \todo{Add \texttt{getLine()} method}

\chapter{Python Reference} \label{chap:python}

Python's built-in support for arbitrary-size integers (using
\texttt{BigInteger} in Java is a pain!) and built-in dictionaries with
lightweight syntax make it attractive for certain kinds of problems.

\section{Template}

Below is a basic template showing how to read typical contest problem
input in Python:

\pycode{python/template.py}

\todo{Mention basic Python data structures such as set, deque, list methods}

\chapter{Data Structures}

\section{Pair}

Java has a \texttt{Pair} class in \texttt{javafx.util.Pair}, but one
can't necessarily assume that this will be available in a contest
environment.  The following simple class suffices to store pairs of
values, which is especially useful in representing \eg coordinates in
a 2D grid.  Note in a competitive programming context we don't bother
adding getter and setter methods; the \texttt{a} and
\texttt{b} fields are public so we can just access them directly.

\code{data-structures/Pair.java}

Below is a variant which implements the \texttt{Comparable}
interface. \texttt{CPair} objects sort in lexicographic order: first
by the first component, and then ties in the first component are
broken by the second component.

\code{data-structures/CPair.java}

\section{Bag}

\kattis{cookieselection, kattissquest}

A \emph{bag} is a collection of elements where order does not matter
(like a set) but multiplicity does matter, \ie there can be duplicates
and we need to keep track of how many duplicates there are of each
item.  Bags are not needed often but can occasionally be useful.  It
is not too hard to build a bag as a map from items to integer counts,
but there are a few corner cases so it's worth copying a well-tested
implementation instead of writing one from scratch.

The implementation below is based on a \javadoc{util}{TreeMap}
\link{sec:map}, and hence supports operations like \texttt{first()}
and \texttt{last()}.  If desired one could easily change the
\texttt{TreeMap} to a \texttt{HashMap} and remove the methods which
are no longer supported, although the factor of $O(\lg n)$ is unlikely
to make a practical difference.

\code{data-structures/TreeBag.java}

\section{Union-find}

\kattis{forestfires, kastenlauf, ladice, numbersetseasy, unionfind, virtualfriends,
  watersheds, wheresmyinternet}

A union-find structure can be used to keep track of a collection of
disjoint sets, with the ability to quickly test whether two items are
in the same set, and to quickly union two given sets into one.  It is
used in Kruskal's Minimum Spanning Tree algorithm \link{sec:kruskal},
and can also be useful on its own (see the above Kattis problems for
examples).  \texttt{find} and \texttt{union} both take essentially
constant amortized time.

\code{data-structures/UnionFind.java}

The above code can easily be enhanced to keep track of the number of
sets (initialize to \texttt{n}; subtract one every time \texttt{union}
hits the \texttt{ru != rv} case), or to keep track of the actual size
of each set instead of just the rank/height (keep a size for each
index; initialize all to 1; add sizes appropriately when doing
\texttt{union}).

\section{Tries}

\kattis{boggle, heritage, herkabe, phonelist}

The code below is a very simple implementation of a trie---there are
many other methods that could be added, and it is not very efficient
since it repeatedly uses the $O(n)$ \texttt{substring} operation as it
recurses down the trie, but it is sufficient for some problems.

\todo{More efficient/full-featured Trie class}

\code{data-structures/Trie.java}

Tries are intimately connected with MSD radix sort, which can be
thought of as equivalent to building a trie and then traversing it in
order.  However, no implementation of radix sort actually builds an
intermediate trie.  Sometimes it can be helpful to think about a
problem in terms of a trie, but never actually implement/materialize
the trie at all (\inlinekattis{herkabe}): just do a modified radix
sort, first grouping strings by their first character, then recursing
on each group, keeping track of needed auxiliary information
(\eg depth) along the way.

\section{Adjustable priority queue}
\label{sec:adj-pq}

\kattis{flowerytrails}

As discussed in \fulllink{sec:pq}, Java's \texttt{PriorityQueue} class
has no way to efficiently alter the priority of an item already stored
in the queue; simply removing and re-adding the item does the trick
but takes $O(n)$ time.  The efficiency of this operation really does
make a difference in the asymptotic performance of Dijkstra's
algorithm \link{sec:dijkstra}, and occasionally it really needs to be
$O(\lg n)$ in order to meet the time limits (\emph{e.g.}
\inlinekattis{flowerytrails}). A suitable implementation of a priority
queue with $O(\lg n)$ priority adjustment is shown below.  The key
idea is to keep a hash table on the side which can be used to quickly
find the index of any item stored in the priority queue; of course,
the hash table has to be kept suitably updated whenever items are
shuffled in the heap.  The \texttt{adjust(e)} method is used to inform
the priority queue that the priority of item \texttt{e} has changed,
so that the queue has an opportunity to move the item if necessary to
reestablish the heap invariants.

\code{data-structures/AdjustablePQ.java}

\section{Segment trees and Fenwick trees}

See \fulllink{sec:range-queries}.

\chapter{Search}

\section{Complete search}

\kattis{bing, classpicture, coloring, cycleseasy, dancerecital,
  lektira, freefood, gepetto, kastenlauf, mjehuric, paintings, prozor,
  rectanglesurrounding, reducedidnumbers, reseto, sheldon, shuffling,
  weakvertices, wheels, transportationplanning}

See CP3 for a fuller discussion of complete search, aka brute force,
and a list of relevant techniques (nested loops, recursive
backtracking, \emph{etc.}).  Just remember that there's no need to
code anything more sophisticated if a back-of-the-envelope analysis
shows that a simple complete search will finish under the time
limit. (Although some kinds of complete search can themselves be
rather sophisticated.  For example, see \fulllink{chap:bittricks}.
Some of the above problems are much harder than others!)

Sometimes complete search isn't in and of itself the full solution to
a problem, but the problem is set up so that a subpart can be done via
complete search, to keep the solution complexity from getting out of
hand and allowing you to focus your efforts on the more
``interesting'' part of the problem.

\section{Binary search}
\label{sec:binary-search}

If you need to do a traditional binary search---that is, finding the
index where a given element occurs in a sorted array---you should just
use the standard \texttt{Arrays.binarySearch} method.  However, the
underlying idea of binary search applies in many more contexts.

\subsection*{Binary search on a real interval}
\label{sec:binary-search-real}

\kattis{bottles, cheese, speed, suspensionbridges, tetration}

This is probably the most common form of binary search in competitive
programming.  Given a function $f$ which is monotonic (\emph{i.e.}
always increasing, or always decreasing) on a given interval of the
\emph{real} line $[a, b]$, find $a \leq x \leq b$ such that $f(x)$ is
equal to some target value.  This can be accomplished by
straightforward binary search: keep track of a current subinterval
$[x_L, x_H]$; at each step, evaluate $f$ at the midpoint
$m = (x_L + x_H)/2$ of the interval, and update $x_L$ or $x_H$ to $m$
depending on whether the value of $f$ is too small or too big,
respectively.  Iterate until $x_H - x_L$ is within an appropriate
tolerance (or simply iterate a fixed number of times---$50$ should be
plenty), and return $(x_L+x_H)/2$.  This is actually easier than
traditional binary search since one doesn't have to worry about
indexing, off-by-one errors, and the like.

The main trick is to realize when this technique is applicable.
Sometimes the function $f$ is plainly stated in the problem
description, but sometimes the thing being searched for is more
subtle.  Whenever a problem asks for a floating-point number as the
answer, it's worth considering whether you can binary search for it.

\todo{Example code?}

\subsection*{Binary search on an integer interval}
\label{sec:unbounded-binary-search}

\kattis{guess, eko, freeweights, inversefactorial, reservoir}

Suppose we again have a function $f$ and want to find a value $n$ such
that $f(n)$ is equal to some target value $t$---except that $f$ is defined
on the \emph{integers} instead of the real numbers.  We can again use
binary search---but we have to be much more careful about potential
off-by-one errors.

\newcommand{\vlo}{\mathit{lo}}
\newcommand{\vhi}{\mathit{hi}}
\newcommand{\vmid}{\mathit{mid}}

\begin{itemize}
\item In the basic version, we simply want to find $n$ such that
  $f(n) = t$, or report that no such $n$ exists.  In this case it
  works well to use a half-open interval, that is, we maintain the
  invariant that possible values of $n$ lie in the interval
  $[\vlo, \vhi)$, including $\vlo$ but \textbf{excluding}
  $\vhi$.  This has the advantage that the size of the
  remaining interval can be computed as simply $\vhi - \vlo$,
  and an appropriate condition for the loop is $\vhi - \vlo > 0$.

  The midpoint of $[\vlo, \vhi)$ can be computed as
  $\vmid = (\vlo + \vhi)/2$;\footnote{Or
    $\vmid = \vlo + (\vhi - \vlo)/2$, if you are worried about
    $\vlo + \vhi$ overflowing, but this is unlikely to ever be an
    issue in a competitive programming context.} $\vmid$ always
  lies within the interval, even if $\vhi - \vlo = 1$ (the rounding
  behavior of integer division plays a crucial role).

  If $\vmid$ does not hold the target, one must then either update
  $\vhi$ to $\vmid$, or $\vlo$ to $\vmid + 1$ (not $\vmid$!) depending
  on whether the item at $\vmid$ is larger or smaller than the target,
  respectively.

  \todo{Example code}

\item A slightly more sophisticated variant is where we need to find
  the largest $n$ such that $f(n) \leq t$, or the smallest such that
  $f(n) \geq t$, or something similar. This requires even more care.
  In this situation it tends to be better to use a closed interval
  $[\vlo, \vhi]$, and using great care to update $\vlo$ and $\vhi$
  appropriately (to $\vmid - 1$, $\vmid$, or $\vmid+1$) depending on
  the desired properties of the value being searched for.

  In this scenario, when there can be duplicate values of $f(n)$, it's
  not possible to stop the search early, since any given $n$ for which
  $f(n) = t$ may not be the optimal one.  One must continue searching
  until the interval has reached size $1$, and then return the sole
  remaining value.

  \todo{Example code}
\end{itemize}

\subsection*{Unbounded binary search}

Consider the following problem: given an increasing function $f$ and a
target value $t$, find the smallest positive value of $x$ such that
$f(x) \geq t$. (The domain of $f$ can be either the reals or the
integers.)

\todo{Find some example Kattis problems that need an initial unbounded
  search?}

\section{Ternary search}

\kattis{brocard, euclideantsp, infiniteslides, janitortroubles, dailydivision}

Ternary search can be used to find the minimum or maximum of a
function which is concave or convex on a given interval (that is, the
function only decreases until the minimum and then only increases, or
vice versa).  Binary search does not apply in this case, since just by
looking at the value of the function at the midpoint of the interval,
it is impossible to know whether we should recurse on the left or
right half of the interval.

Suppose we are currently considering the interval $[L,R]$ and looking
for the minimum of a function $f$ on the interval.  We compute the two
points $1/3$ and $2/3$ of the way through the interval, namely
$m_L = (2L + R)/3$ and $m_R = (L + 2R)/3$.
\begin{itemize}
\item If $m_L < m_R$, then we know the minimum can't be to the right
  of $m_R$ (because then it would increase from $m_L$ to $m_R$ and
  then decrease---but we assume the function decreases until the minimum
  and then only increases after that).  Hence, we can recurse on the
  interval $[L, m_R]$.
\item If $m_L > m_R$, we can likewise recurse on $[m_L, R]$.
\item If $m_L = m_R$, we can recurse on $[m_L, m_R]$ (though lumping
  this case in with either of the above two cases works fine and
  requires writing less code).
\end{itemize}
In any case, we decrease the size of the interval by at least $1/3$
with each iteration, so we need only a logarithmic number of
iterations relative to the ratio between the starting interval size
and the desired accuracy.

\todo{Example code}

\subsection*{Integer ternary search}

When ternary searching over an interval of \emph{integers}, much of
the same advice applies as for binary search (see the previous
section).  However, care must be taken with the stopping conditions;
depending on exactly how the recursion works it is possible to end up
in a scenario where it loops infinitely on an interval of size 1 or 2.
Even if it is possible to come up with an elegant design that does
not require any special cases, it may be easiest to simply stop the
loop when the interval has size 2 or smaller, and then simply check
the few remaining items manually.

\chapter{Graphs}

\section{Graph basics} \label{sec:graph-basics}

\todo{Directed, undirected, weighted, unweighted, self loops, multiple edges}
\todo{characterization of trees}
\todo{New virtual source/sink node trick}

\kattis{chopwood}

\section{Graph representation}

\todo{Adjacency matrix, adjacency maps.  Edge objects. Implicit
  graphs.}

Figure~\ref{fig:horrorlist} has a sample solution for
\inlinekattis{horrorlist} which builds an adjacency map
representation of an undirected graph.

\todo{State space search with complex states: make a class, implement
  Comparable, use TreeMap}

\section{Breadth-First Search}
\label{sec:bfs}

\kattis{ballsandneedles, brexit, buttonbashing, collapse, erdosnumbers, grapevine, horrorlist, mazemakers}

Breadth-first search (BFS) can be used to find single-source shortest
paths (\emph{i.e.} shortest paths from a particular starting vertex to
all other vertices) in an unweighted graph. BFS comes up often in many
different guises, so it's worth being very familiar with BFS and its
variants. Below is pseudocode showing a generic BFS implementation.
Important invariants:
\begin{itemize}
\item Every vertex in $Q$ has already been marked visited. (This is
  important since it prevents vertices from being added to $Q$
  multiple times.)
\item $Q$ only contains vertices from at most two (consecutive) levels
  at a time.
\end{itemize}

\begin{center}
\begin{minipage}{0.8\textwidth}
\begin{algorithm}[H]
  \begin{algorithmic}[1]
\State $s \gets$ starting vertex
\State Mark $s$ visited
\State $Q \gets$ new queue containing only $s$
\State $\mathit{level}[s] \gets 0$
\While{$Q$ not empty}
  \State $u \gets Q.\mathit{remove}$
  \For{each neighbor $v$ of $u$}
    \If{$v$ is not visited}
      \State $\mathit{level}[v] \gets \mathit{level}[u] + 1$ \Comment{Optionally mark level}
      \State Add $v$ to $Q$
      \State Mark $v$ visited
      \State $\mathit{parent}[v] \gets u$ \Comment{Optionally record parent}
    \EndIf
  \EndFor
\EndWhile
  \end{algorithmic}
\caption{BFS}
\end{algorithm}
\end{minipage}
\end{center}

Some options/variants:
\begin{itemize}
\item The $\mathit{level}$ array shown above is optional, and can be
  omitted if not needed.  Sometimes it makes sense to have the
  $\mathit{level}$ array do double-duty to also track visited vertices:
  if the $\mathit{level}$ of every vertex is initialized to some
  nonsensical value such as $-1$ or $\infty$, then a vertex is visited
  iff its $\mathit{level}$ is not equal to the initial value.

  Figure~\ref{fig:horrorlist} shows a sample solution for
  \inlinekattis{horrorlist}, exhibiting a BFS with level labelling.
\item The parent map is also optional, and can be used to reconstruct
  an actual shortest path from $s$ to any vertex, by starting with the
  end vertex and iteratively following parents backwards until
  reaching $s$.
\item If you want to compute shortest paths from \emph{any} of a set
  of starting vertices, simply replace the initialization of $s$ with
  the desired set (\emph{i.e.} mark them all visited, add them all to
  $Q$, and set their $\mathit{level}$ to $0$ before starting the loop;
  the loop itself does not change) (\inlinekattis{zoning}).
\item Replacing $Q$ with a stack results in a depth-first rather than
  breadth-first search (although often it makes more sense to
  implement a DFS recursively; see \link{sec:dfs}).
\end{itemize}

\begin{figure}
  \small
  \code{graph/horrorList.java}
  \caption{Sample solution for \texttt{horrorlist} (Adjacency set
    representation; BFS with level labelling)}  \label{fig:horrorlist}
\end{figure}

\todo{Applications of BFS: identify reachable vertices; identify
  (weakly) connected components; identify bipartite graphs/odd cycles
  (detect cross-edges with map of level sets)}

\section{DFS and SCCs} \label{sec:dfs}

\todo{Code for DFS, start/finish labelling, top sorting, Tarjan's SCC algorithm}

\section{Topological sorting} \label{sec:topsort}

\kattis{builddeps, easyascab, eatingeverything, excavatorexpedition,
  mravi, promotions, reactivity, runningmom, succession}

% , ccsc18.topicx

A \emph{topological sort} of a directed graph $G$ is a list of
vertices such that whenever there is an edge from $u$ to $v$, $u$
comes before $v$ in the list; $G$ has a topological sort if and only
if it is acyclic.  Topological sorting can thus be used to detect the
presence of cycles. It is also often used in conjunction with dynamic
programming (\inlinekattis{eatingeverything, excavatorexpedition,
  mravi}): if we need to compute some value each vertex
such that the value can be computed once we already know the values
for all the outgoing (or incoming) neighbors, topological sort gives
us the right order for computing the values.

There are two main methods to do a topological sort.  Method 1 (Kahn's
Algorithm) is to repeatedly remove nodes with no incoming edges (or
dually, nodes with no outgoing edges).  Empirically this seems to be
faster than Method 2, but is perhaps a bit more code.  Pseudocode is
as follows:

\begin{center}
\begin{minipage}{0.8\textwidth}
\begin{algorithm}[H]
  \begin{algorithmic}[1]
    \Require Directed graph $G = (V,E)$.
    \State $T \gets$ empty list (to store topsort)
    \State $Z \gets$ empty queue (to store nodes with $0$ indegree)
    \State $\mathit{in} \gets$ dictionary mapping all vertices to
    their indegree
    \State Put all vertices with indegree $0$ into $Z$
    \While{$Z$ is not empty}
      \State $v \gets Z.\mathit{dequeue}$
      \State append $v$ to $T$
      \For{each $u$ adjacent to $v$}
        \State decrement $\mathit{in}[u]$
        \If{$\mathit{in}[u] = 0$}
          \State add $u$ to $Z$
        \EndIf
      \EndFor
    \EndWhile
  \end{algorithmic}
  \caption{{\sc TopSort}(G)}
  \label{alg:topsort}
\end{algorithm}
\end{minipage}
\end{center}

If the queue becomes empty before all vertices have been added to the
topsort, then a cycle exists.

For a sample implementation of this algorithm, see the solution to
\inlinekattis{succession} at \url{https://github.com/Hendrix-CS/programming-team/blob/master/solved/Succession.java}.

The second method is to do a recursive DFS: simply add each vertex to
a list just \emph{after} recursively processing all its neighbors;
this yields a topsort in reverse order.

\begin{center}
  \begin{minipage}{0.8\textwidth}
    \begin{algorithm}[H]
      \begin{algorithmic}[1]
        \Function{TopSort-DFS}{G}
        \State $T \gets$ empty list/stack to hold topsort
        \ForAll{$v \in V$}
        \If{$v$ is not visited}
        \Call{DFS}{$v$, $T$}
        \EndIf
        \EndFor
        \State \Return{$T$}
        \EndFunction \\

        \Function{DFS}{$x$, $T$}
        \State Mark $x$ visited
        \ForAll{$(x,y) \in E$}
          \If{$y$ is not visited}
            \Call{DFS}{$y$, $T$}
          \EndIf
        \EndFor
        \State Add $x$ to $T$
        \EndFunction
      \end{algorithmic}
      \caption{Topological sort via DFS}
      \label{alg:topsort-dfs}
    \end{algorithm}
  \end{minipage}
\end{center}

\section{Single-source shortest paths (Dijkstra)} \label{sec:dijkstra}

\kattis{bigtruck, blockcrusher, coffeedate, george, getshorty, kitchen,
  shortestpath1, shortestpath2, showroom, walkway}

Dijkstra's algorithm is the standard algorithm for solving the
\emph{single-source shortest path} problem in weighted, directed
graphs.  That is, given a graph with (possibly) directed edges and a
weight on each edge, Dijkstra's algorithm can find the shortest
directed path from a single chosen start vertex to every other vertex
in the graph (where the length of a path is the sum of the weights on
the edges). If you want to find the shortest path in a weighted,
undirected graph, just make a directed graph with edges going both
directions between each pair of vertices.  Figure~\ref{fig:dijkstra}
has a basic implementation.
\begin{figure}
  \small
  \code{graph/Dijkstra.java}
  \caption{$O(VE)$ Dijkstra's algorithm} \label{fig:dijkstra}
\end{figure}
\begin{warning}
  Since Java's \javadoc{util}{PriorityQueue} class does not have a
  ``decrease key'' method, on line 28 we have to instead do a
  \texttt{remove} followed by an \texttt{add}; but \texttt{remove} is
  $O(n)$, making the whole algorithm $O(VE)$.  If you really need
  $O(E \lg V)$ performance (\inlinekattis{flowerytrails}), you can use
  an \fulllink{sec:adj-pq}.  In some situations you can also simply
  call \texttt{add} without calling \texttt{remove}; see the
  discussion below.

  \todo{using adjust, you have to decide whether to 'add' or 'adjust'}
\end{warning}

There are many possible variants of this basic template; here are a
few.

\begin{itemize}
\item The given code explores the \emph{entire} graph.  However, if
  you have a particular target vertex in mind you can stop early once
  you find it: just \texttt{break} out of the loop if removing the
  next node from the priority queue yields the target node, since at
  that point we are guaranteed that we know the shortest path from the
  start node to the target node.
\item If the vertices of your graph are not naturally represented as
  integers in the range $0 \dots n-1$, one could modify the algorithm
  to use \texttt{Map}s in place of the \texttt{parent} and
  \texttt{dist} arrays.  Alternatively, it may be easier to deal with
  this outside of Dijkstra's algorithm: just arbitrarily assign
  indices to vertices and use a \texttt{Map} or two to keep track of
  the assignment.  Then run Dijkstra using the assigned vertex indices
  and translate the result back to the original vertices.
\item If the priority queue contains objects whose priority never
  changes once they are put in the priority queue (note that the
  example code in Figure~\ref{fig:dijkstra} does \emph{not} have this
  property, since \texttt{Integer}s in the PQ are compared by the
  value stored in the external array \texttt{dist}, which can change)
  then it can be an optimization to simply call \texttt{pq.add(next)}
  without calling \texttt{pq.remove(next)} first. The priority queue
  will end up with multiple copies of the same node, each with a
  different priority, but this is not a problem; when removing the
  next node from the PQ just ignore it if it has already been visited.
  (\inlinekattis{nikola})
\item Dijkstra's algorithm uses addition to combine the weights of
  consecutive edges and \texttt{min} to pick the shortest path among
  parallel options.  However, there are other pairs of operations one
  can use with the same basic algorithm template.\footnote{The details
  of which properties of the operations are needed for this to work
  are too far outside the scope of this document; see \todo{XXX}}
  \begin{itemize}
  \item Using \texttt{min}/\texttt{max} in place of
    \texttt{+}/\texttt{min} yields an algorithm which finds the path
    with the maximum possible minimum weight (\inlinekattis{vuk,
      crowdcontrol, muddyhike}).  For example, if the edge weights are
    thought of as capacities, and the capacity of a path is equal to
    the minimum capacity of any of its edges (\emph{i.e.} the
    bottleneck) then this corresponds to finding maximum-capacity
    paths.  One must be careful to:
    \begin{itemize}
    \item update the comparison operation for the priority queue to
      use \texttt{min} instead of \texttt{max} (\emph{e.g.} by
      switching to \texttt{dist[v] - dist[u]} instead of
      \texttt{dist[u] - dist[v]}),
    \item initialize all the entries of \texttt{dist} to an
      appropriate identity value for \texttt{max} such as \texttt{0},
      \texttt{-1}, or \texttt{-INF} instead of \texttt{INF},
    \item change the definition of \texttt{nextDist} to use
      \texttt{min} instead of \texttt{+}, and
    \item change the comparison of \texttt{nextDist} and
      \texttt{dist[next]} to use \texttt{>} instead of \texttt{<}.
    \end{itemize}
  \item If we have a directed graph with edge weights corresponding to
    probabilities, where the probability of a path is defined as the
    product of the probabilities of its edges, then Dijkstra's algorithm
    with \texttt{*}/\texttt{max} finds highest-probability paths.
    Similar modifications have to be made as in the previous example.
  \end{itemize}
\item One can modify the basic algorithm to keep track of extra
  information, such as the \emph{number} of shortest paths from the
  start to any given node: add to the count when finding a new path
  equal in weight to the previous best-known path; reset the count
  when finding a shorter path than previously known
  (\inlinekattis{visualgo}).
\item Dijkstra can also deal with edges whose weight depends on the
  time they are reached (think of \emph{e.g.} bus routes, where you
  may have to wait a while for the next bus to come depending on what
  time you reach the stop). (\inlinekattis{coffeedate})
\end{itemize}

\section{All-pairs shortest paths (Floyd-Warshall)} \label{sec:floydwarshall}

\kattis{crosscountry, allpairspath, shoppingmalls, transportationplanning, units}

The \emph{all-pairs shortest path} problem is to find the shortest
path in a (directed, weighted) graph between \emph{any} pair of
vertices.  Typically the idea is to precompute some table(s)
and then be able to quickly look up any pair of vertices to find the
distance between them.  This could be done by running \emph{e.g.}
Dijkstra's algorithm once from every vertex, but that takes at least
$O(VE \log V)$ (which is $O(V^3 \log V)$ for a dense graph) and
doesn't work if there are negative edge weights. The Floyd-Warshall
algorithm runs in $O(V^3)$ no matter how many edges there, can
handle negative edge weights, and is just a few lines of code.

Note this only works when $|V|$ is small enough for a cubic algorithm
to fit in the time limits, typically something like $|V| \leq 400$,
though I have seen examples with $|V|$ even up to $1000$ that work.
Each individual loop of Floyd-Warshall is only a few operations so the
constant factor is very small.

Assume the vertices in $G$ are labelled $\{0, \dots, n-1\}$.  Create a
2D matrix of distances $d$ and initialize it like so:
\[ d[i][j] =
  \begin{cases}
    0      & \text{if $i = j$} \\
    w_{ij} & \text{if there is an edge $i \to j$} \\
    \infty & \text{otherwise}
  \end{cases}
\]
If there can be multiple edges from $i$ to $j$, be sure to set
$d[i][j]$ to the \emph{minimum} of all the edge weights. In practice,
for $\infty$, just use a value that is much larger than any other
values that could occur in the problem.

Then the Floyd-Warshall algorithm is as follows.  We iterate $k$ from
$0$ to $n-1$; $k$ represents the intermediate vertex we will consider.
Then for every possible pair of vertices $u$ and $v$, we check if
there is a way to get from $u$ to $k$ and a way to get from $k$ to
$v$, and the sum of these distances is less than the current shortest
distance from $u$ to $v$.  If so, we update it.

\begin{minted}{java}
for (int k = 0; k < n; k++)
  for (int u = 0; u < n; u++)
    for (int v = 0; v < n; v++)
      if (d[u][k] < INF && d[k][v] < INF)
        d[u][v] = Math.min(d[u][v], d[u][k] + d[k][v]);
\end{minted}

After running the above loops, $d[i][j]$ will contain the length of
the shortest path from $i$ to $j$.  If there is no path from $i$ to
$j$ then the length will be $\infty$.  Analyzing the running time of
this algorithm is a Data Structures student's dream: there are
literally three nested loops which each iterate $|V|$ times, so the
running time is $O(V^3)$.

Variants:

\begin{itemize}
\item If you want to detect the presence of negative cycles, you can
  use the fact that after running the main algorithm, $d[i][i] < 0$ if
  and only if vertex $i$ is contained in some negative cycle.  Hence
  if there are no negative numbers along the diagonal of the matrix,
  the graph is negative-cycle-free.  However, if there are
  negative cycles then the distances found may not be valid (if there
  is a negative cycle along a path from $u$ to $v$ then one could travel
  around the cycle any number of times before finally going to $v$).

\item If you actually want to distinguish negative-cycle-free paths
  (for which the computed minimum distance is valid) from others, you
  can run the following additional code, which propagates information
  about negative cycles:

  \begin{minted}{java}
    for (int u = 0; u < n; u++)
      for (int v = 0; v < n; v++)
        for (int k = 0; d[u][v] != -INF && k < n; k++)
          if (d[u][k] < INF && d[k][k] < 0 && d[k][v] < INF)
            d[u][v] = -INF;
  \end{minted}

  For each pair of vertices $u, v$, we check all possible intermediate
  nodes $k$. If there is a path $u \to k$ and a path $k \to v$,
  and $k$ is part of some negative cycle, then we set the distance
  from $u$ to $v$ to $-\infty$ to signify that the length of a path
  from $u$ to $v$ can be arbitrarily small.

\item Sometimes you want to know not only the \emph{length} of the
  shortest path, but the actual shortest path itself.  This can be
  accomplished by keeping a 2D array $\mathit{next}$ such that
  $\mathit{next}[i][j]$ stores the next vertex along the shortest path
  from $i$ to $j$.  Initialize $\mathit{next}[i][j]$ to $j$ whenever
  there is an edge from $i$ to $j$ (it does not matter what value it
  has otherwise).  Then update the \texttt{if} statement in the inner
  loop as follows:
  \begin{minted}{java}
    if (d[u][k] < INF && d[k][v] < INF && d[u][k] + d[k][v] < d[u][v]) {
      d[u][v] = d[u][k] + d[k][v];
      next[u][v] = next[u][k];
    }
  \end{minted}
  Now after running the algorithm, the shortest path from $u$ to $v$
  can be recovered by looking up $u_2 = \mathit{next}[u][v]$, then
  $u_3 = \mathit{next}[u_2][v]$, and so on, until $v$ is reached.

\end{itemize}

\section{Min spanning trees (Kruskal)} \label{sec:kruskal}

\kattis{drivingrange, islandhopping, jurassicjigsaw, lostmap, minspantree, treehouses}

\todo{Create a union-find; sort all the edges by size; for each edge
  smallest to largest, pick it and connect its endpoints if they
  aren't already connected.}

\section{Eulerian paths} \label{sec:eulerian-path}

\section{Max flow/min cut}

\kattis{copsandrobbers,escapeplan,gopher2,guardianofdecency,marblestree,maxflow,mincut,paintball,pianolessons,waif}

A \emph{flow network} is a directed, weighted graph where the edge
weights (typically integers) are thought of as representing
\emph{capacities} (\eg imagine pipes of varying sizes).  The \emph{max
  flow problem} is to determine, given a flow network, the maximum
possible amount of \emph{flow} which can move through the network
between given source and sink vertices, subject to the constraints
that the flow on any edge is no greater than the capacity, and the sum
of incoming flows equals outgoing flows at every vertex other than the
source or sink.  Flow networks can be used to model a wide variety of
problems.

\todo{Enumerate a few problem types: item assignment; max bipartite
  matching; min cut}

\todo{choose directed/undirected edges carefully!}

\todo{Requires vertices $0 \dots n-1$: either carefully keep track of
  which numbers are for which vertices, or use lookup tables}

Dinitz' Algorithm is probably the best all-around algorithm to use for
solving max flow problems in competitive programming.  It takes
$O(V^2 E)$ in theory (although is often much faster in practice).  In
the special case where we are modelling a bipartite matching problem,
Dinitz' Algorithm reduces to the Hopcroft-Karp algorithm which runs in
$O(E \sqrt{V})$.

\code{flow/Dinitz.java}

\todo{Include a sample solution using a flow network}

\todo{Variants: Multiple sources/sinks? Use trick of adding a new
  source/sink with infinite capacity edges. Vertex capacities?
  Turn each vertex into a new edge.}

\chapter{Dynamic Programming}
\label{chap:dyn-prog}

\kattis{balanceddiet, drivinglanes, walkforest}

\todo{subset sum}
\todo{knapsack, longest common subsequence}
\todo{longest increasing subsequence ($O(n^2)$ and $O(n \lg n)$, see \url{https://stackoverflow.com/questions/2631726/how-to-determine-the-longest-increasing-subsequence-using-dynamic-programming})}

\chapter{Sequences and strings}

\section{Longest Increasing Subsequence (LIS)}

\kattis{increasingsubsequence,longincsubseq,manhattanmornings,signals}

A \emph{subsequence} of a sequence is a subset of the elements, taken in
order, but not necessarily contiguous.  (By contrast, a contiguous
subset of elements is often referred to as a \emph{subinterval}.)  For
example, $[1, 3, 4, 7]$ is a subsequence of
$[1, 2, 3, 4, 5, 6, 7, 8]$.  Given a sequence of integers (or any
elements which can be ordered), the \emph{longest increasing
  subsequence} (LIS) problem is to find the longest subsequence which
is in strictly increasing order.  For example, given the sequence
$[9,2,8,10,5,4,20,16,7,1]$, one increasing subsequence is
$[2, 5, 16]$, but it is not the longest.  There are several increasing
subsequences of length $4$, such as $[2,8,10,16]$, and it turns out
that this is the longest possible.

Conceptually, to compute the LIS of a sequence, we first build a
\emph{downravel}, a set of nonincreasing subsequences which partition
the original sequence.  We keep these subsequences as a list of
stacks, and maintain the invariant that their top elements are always
sorted from smallest to biggest.  We iterate through the elements of
$S$ and push each onto the leftmost possible stack whose top is
$\geq$ the element being added.

\begin{center}
  \begin{minipage}{0.8\textwidth}
    \begin{algorithm}[H]
      \begin{algorithmic}[1]
        \Function{Downravel}{S}
        \State $D \gets$ empty list of stacks
        \ForAll{$x \in S$}
          \State $k \gets$ first stack in $D$ whose top is $\geq x$
          \If{no such $k$ exists}
            \State Add a new singleton stack containing $x$ to the end
            of $D$
          \Else
            \State Push $x$ onto $k$
          \EndIf
          \EndFor
        \State \Return{$D$}
        \EndFunction
      \end{algorithmic}
      \caption{Building a downravel of a sequence $S$}
      \label{alg:downravel}
    \end{algorithm}
  \end{minipage}
\end{center}

\todo{add some pictures?}

The length of the LIS is the same as the length of the downravel $D$.
Na\"ively, this runs in $O(n^2)$ time, since for each of the the $n$
elements in $S$, we have to search through up to $O(n)$ stacks in $D$
to find the right one to push.  However, there are several possible
optimizations.

\begin{itemize}
\item First, although the stacks are conceptually helpful, we do not
  actually need to store them. It's enough to simply store the current
  top element of each stack.  So instead of having a list of stacks we
  just have a list of elements.
\item Second, since this list will always be sorted from smallest to
  biggest, we can use a binary search to find the proper place to push
  each new element.  This brings the running time down to a very
  respectable $O(n \lg n)$.
\end{itemize}

Of course, we might want not just the \emph{length} of the LIS but an
actual LIS itself.

\begin{itemize}
\item First of all, we need to modify $D$ so it stores not the elements
  themselves but their indices in $S$.  This requires a bit of extra
  indirection while doing a binary search for the correct place to put
  each new element.
\item We keep an extra array $\mathit{back}$ such that
  $\mathit{back}[i]$ stores the index of an element that could come
  before $S[i]$ in a LIS up to and including $S[i]$.  Every time we
  put a new element $S[i]$ into $D$, we set $\mathit{back[i]}$ to the
  previous value in $D$---that is, the index of the element currently
  on top of the previous stack in the list.
\item After running the algorithm we start with the item represented by
  the last entry in $D$, and then keep following $\mathit{back}$ links
  to get each previous item.  This yields a LIS in reverse order.
\end{itemize}

\todo{include some examples and pictures}

\code{sequence/LISCompact.java}

\section{Longest Common Subsequence (LCS) via LIS}
\kattis{inflagrantedelicto, princeandprincess}

Given two sequences, the \emph{longest common subsequence} (LCS)
problem is to find the longest sequence which is a subsequence of
both. This comes up in quite a few real-world applications including
DNA processing and diffing (\emph{i.e.} what Github does when it shows you
which lines have changed).

Most generally, the LCS of two sequences with lengths $m$ and $n$ can
be computed via \fulllink{chap:dyn-prog} in $O(mn)$ time.  However, in
the special case that the sequences do not have too many repeated
elements, it is possible to solve it more quickly, as follows.

Suppose the sequences are called $A$ and $B$.
\begin{itemize}
\item Construct all pairs of indices $(i,j)$ such that $A[i] = B[j]$,
  that is, all pairs of locations where $A$ and $B$ agree.
\item Sort these pairs lexicographically, that is, sort them first by
  $i$ and break ties by $j$.
\item Now make an array consisting of just the $j$ values from these
  sorted pairs.
\item A longest increasing subsequence in this array of $j$ values
  gives the length of a LCS of $A$ and $B$.  (Exercise for the reader:
  why does this work?)
\end{itemize}

\section{Z-algorithm}

\section{Suffix arrays}

\chapter{Mathematics}

\section{GCD/Euclidean Algorithm} \label{sec:euclid}

The \emph{Euclidean algorithm} can be used to compute the greatest
common divisor of two \textbf{nonnegative} integers. (If you need it
to work for negative numbers as well, just take absolute values
first.)  It runs in logarithmic time.  The \emph{extended Euclidean
  algorithm} not only finds the GCD $g$ of $a$ and $b$, but also finds
integers $x$ and $y$ such that $ax + by = g$.

\kattis{fairwarning, jughard, kutevi, candydistribution}

\code{math/GCD.java}

\section{Rational numbers}

\kattis{bikegears,jointattack,prosjek,prsteni,rationalarithmetic,wheels,zipfsong}

Occasional problems may require dealing with explicit rational values
rather than using floating-point approximations.  If a problem
involves non-integer values but requires being able to test values for
equality \emph{exactly}, then likely rational numbers are required.
The below code for a \texttt{Rational} class is not difficult but it's nice
to have it as a reference. Of course in a real contest situation you
may not need all the methods.

\newpage

\code{math/Rational.java}

\section{Modular arithmetic}

\kattis{crackingrsa,modulararithmetic,pseudoprime,reducedidnumbers}

\begin{warning}
  Java's mod operator \texttt{\%} behaves strangely on negative
  numbers. In many other languages (\eg Python, Haskell) \texttt{a \%
    b} always returns a result between $0$ and $b-1$; however, in Java
  (as in C/C++), if \texttt{a} is negative then \texttt{a \% b} will
  also be negative.  Try adding \texttt{b} first if you need a
  nonnegative result.
\end{warning}

For example, suppose \texttt{i} is an index into an array of length
\texttt{n} and you need to shift by an offset \texttt{o}, wrapping
around in case the index goes off the end of the array.  The obvious
way to write this would be
\begin{minted}{java}
i = (i + o) % n;
\end{minted}
however, this is \textbf{incorrect if \texttt{o} could be negative!}
If we assume that \texttt{o} will never be larger in absolute value
than \texttt{n}, then we could write this correctly as
\begin{minted}{java}
i = (i + o + n) % n;
\end{minted}
If \texttt{o} could be arbitrarily large then we could write
\begin{minted}{java}
i = (((i + o) % n) + n) % n;
\end{minted}
(the first mod operation reduces it to lie between $-n \dots n$;
adding $n$ ensures it is positive; and the final mod reduces it to the
range $[0,n)$).

\subsection*{Modular exponentiation and modular inverses}

Sometimes one needs to compute the modular exponentiation
$b^e \bmod m$ for some base $b$, exponent $e$, and modulus $m$.  Using
repeated squaring, it is possible to do this efficiently even for very
large exponents $e$.  Relatedly, if $b$ is relatively prime to $m$, it
is possible to compute $b^{-1} \bmod m$, the \emph{modular inverse} of
$b$, that is, the unique number $0 < b' < m$ such that
$bb' \equiv 1 \pmod m$.

In Java, probably the easiest way to compute these is using the
\texttt{modPow} method from the \texttt{BigInteger} class
\link{sec:bigint}.  If \texttt{b}, \texttt{e}, and \texttt{m} are
\texttt{BigIntegers}, then \texttt{b.modPow(e, m)} is a
\texttt{BigInteger} that represents $b^e \bmod m$.  The exponent
\texttt{e} can also be negative; in particular, if \texttt{e} is $-1$
then \texttt{b.modPow(e,m)} will compute the inverse of
\texttt{b} modulo \texttt{m}.

It is also useful to know how to compute modular exponentiation and
inverses manually, in case you need some sort of variant version, or
if \texttt{BigInteger} is not fast enough.

\textbf{Modular exponentiation} can be computed by repeated squaring.
The basic idea is to compute $b^e$ by splitting up $e$ into a sum of
powers of two (according to its binary expansion), raising $b$ to each
power of two and taking the product.  This can be done efficiently
since we can get from $b^{2^k}$ to $b^{2^{k+1}}$ just by squaring.

  \begin{warning}
    Even if you need the answer modulo an \texttt{int} value such
    as $10^9 + 7$, it is important to use \texttt{long} in the method
    below: the product of two \texttt{int} values does not necessarily
    fit in an \texttt{int}, even if the very next step will reduce it
    modulo $m$ back into the range of an \texttt{int}.
  \end{warning}

\code[2][10]{math/ModExp.java}

Note this correctly computes $0^0 = 1$.  It would be possible to add a
special case for when $b = 0$ and $e \neq 1$, to avoid multiplying $0$
by itself a bunch of times, but it's hardly worth it.

\textbf{Modular inverses} can be computed using the extended Euclidean
algorithm \link{sec:euclid}.  In particular, suppose $a$ and $b$ are
relatively prime, that is, their GCD is $1$.  In that case the
\texttt{egcd} algorithm will compute numbers $x$ and $y$ such that
$ax + by = 1$.  Taking this equation $(\bmod b)$ yields \[ ax + by
  \equiv ax \equiv 1 \pmod b, \] and so $x$ is the modular inverse of
$a$ modulo $b$ (in practice one may want to reduce $x \bmod b$ so $x$
is between $0$ and $b-1$).

Alternatively, for a prime $p$, Fermat's Little Theorem says that
\[ a^{p-1} \equiv 1 \pmod p \] and hence $a^{p-2}$ is the modular
inverse of $a$ modulo $p$, which can be computed using modular
exponentiation.

\section{Primes and factorization}

Methods for primality testing and prime factorization that may show up
in a contest can be put in two main classes.  First, methods based on
\emph{trial division} are relatively simple to code and work well for
testing just one or a few numbers.  \emph{Sieve} based methods
construct a whole table of primes or factors all at once, and are
often more efficient when many numbers need to be factored or tested
for primality.

\subsection{Trial division}

\kattis{almostperfect,candydivision,crypto,enlarginghashtables,flowergarden,goldbach2,happyprime,iks,listgame,olderbrother,pascal,primalrepresentation}

To test whether a single number is prime, you can use the following
function which performs (somewhat optimized) trial division.  Note
that although there are faster primality testing methods (\eg
Miller-Rabin, Baille-PSW), it is highly unlikely that a contest would
ever require anything more sophisticated than divisibility testing:
Miller-Rabin is not hard to code but it is probabilistic, so a program
using it may give different results on subsequent runs, hardly
suitable for a competitive programming environment; Baille-PSW is
known to be deterministic for numbers up to $2^{64}$, but is much more
complex to code.

Note that \texttt{isPrime} has runtime $O(\sqrt n)$ and is hence
appropriate for numbers up to the maximum size of an \texttt{int}
($\approx 2 \cdot 10^9$); running it on inputs up to the maximum size
of a \texttt{long} is likely to be too slow.

\code[2][10]{math/IsPrime.java}

The following method takes $O(\sqrt n)$ to factor a number into its
prime factorization, also using trial division.  The returned prime
factors will be sorted from smallest to biggest.

\code[4][18]{math/Factor.java}

\subsection{Sieving}

\kattis{industrialspy,nonprimefactors,primereduction,primesieve,reseto}

The term \emph{sieve} comes from the ancient \emph{Sieve of
  Eratosthenes}, a very effective method for generating all the primes
up to a certain bound.  The basic idea is to make a table of all the
numbers from $1$ up to some upper bound $n$ and iterate through the
table. Each time we discover a prime $p$ we ``cross out'' all the
multiples of $p$ in the table; we know a number is prime if it hasn't
yet been crossed out by the time we get to.  This takes time
$O(n \log \log n)$ (essentially linear time) to construct a table for
$1 \dots n$.  The code below uses a \texttt{BitSet} \link{sec:bitset},
which uses less memory than an array of \texttt{boolean}s. Constructing
a \texttt{PrimeSieve} of size $10^8$ should take about a second and
use only about 12 MB of memory; constructing smaller prime sieves
should be quite fast.  Even a \texttt{PrimeSieve} of size
\texttt{Integer.MAX\_VALUE}, \ie $\approx 2 \cdot 10^9$, will fit
quite easily in memory, although constructing it will probably take
too long for most contest problems.  (However, there may be occasional
problems that require building a sieve of this size in order to
precompute some data offline---\ie writing a program that runs for a
few minutes in order to precompute some kind of set or lookup table to
be included in the submitted solution.)

\code{math/PrimeSieve.java}

Instead of simply storing a boolean indicating whether each number is
prime or not, we could also store the smallest prime factor.  We can
still use this to test whether a given number is prime, by checking
whether \texttt{smallest[n] == n}.  But we can also use it to quickly
factor any composite \texttt{n}: simply divide \texttt{n} by
\texttt{smallest[n]} and repeat. We can construct the smallest factor
array using a sieving method similar to \texttt{PrimeSieve}.  The
tradeoff is that this uses much more memory: instead of one bit per
number, we use an entire \texttt{int}, that is, 32 bits.  A
\texttt{FactorSieve} of size $10^8$ will take up around 380 MB.

The \texttt{FactorSieve} class below includes a trivial
\texttt{isPrime} method as well as a \texttt{factor} method, which is
carefully written to work even for \texttt{int} values which are
bigger than the lookup table.

\code{math/FactorSieve.java}

\section{Divisors and Euler's Totient Function}

\kattis{farey,relatives}

\todo{Number of divisors.  Euler's $\varphi$ function: computing
  directly and by sieving.}

\section{Factorial} \label{sec:factorial}

\kattis{eulersnumber, factstone, howmanydigits, lastfactorialdigit,
  inversefactorial, loworderzeros, factovisors}

\todo{Computing factorials; size using logs, etc.  Divisors of factorial.}

\section{Combinatorics} \label{sec:combinatorics}

\kattis{insert, anagramcounting, nine, secretsanta, kingscolors, howmanyzeros}

\todo{Basic principles of combinatorics.  mod $10^9 + 7$. Code for
  computing binomial coefficients. Multinomial coefficients. Large
  binomial coefficients modulo a prime (modular inverse factorial
  tables, Lucas's theorem).}

\begin{warning}
  Remember to use \texttt{long} if you need an answer
  $\bmod (10^9 + 7)$ (which would fit in an \texttt{int}) but
  computing the answer requires \emph{multiplying} $\bmod (10^9 + 7)$.
\end{warning}

\todo{Heap's Algorithm for generating all permutations; next
  permutation.  See Bit Tricks for generating all subsets.}

\todo{PIE?}

\section{Probability}

\todo{Write me}

\chapter{Bit Tricks} \label{chap:bittricks}

\kattis{bits, classpicture, data, flipfive, font, gepetto, hypercube, mazemakers,
  pagelayout, pebblesolitaire, satisfiability, turningtrominos}

\texttt{int} values are represented as a sequence of 32 bits;
\texttt{long} values are 64 bits.  Sometimes it is useful to think
about/work with such values directly as a sequence of bits rather than
as a number. We typically think of the bits as indexed from $0$
starting at the rightmost (least significant) bit.  For example,
\[ \arraycolsep=1pt
  974_{10} =
  \begin{array}{cccccccccc}
    1&1&1&1&0&0&1&1&1&0 \\
      \scriptscriptstyle 9
     &\scriptscriptstyle 8
     &\scriptscriptstyle 7
     &\scriptscriptstyle 6
     &\scriptscriptstyle 5
     &\scriptscriptstyle 4
     &\scriptscriptstyle 3
     &\scriptscriptstyle 2
     &\scriptscriptstyle 1
     &\scriptscriptstyle 0
  \end{array} \]
In general, a $1$ bit at index $i$ has value $2^i$.

One frequently useful point of view is to think of a value of type
\texttt{int}/\texttt{long} as representing a particular subset of a
given set of up to 32/64 items.  The bit at index $i$ indicates
whether item $i$ is included in the subset or not.

Java has built-in operators to manipulate values at the bit level:
\begin{itemize}
\item \texttt{\&} represents bitwise logical AND.  That is, the
  index-$i$ bit of the result is the logical AND of the index-$i$ bits
  of the inputs; each bit index is considered separately.  It is often
  useful to think of \texttt{\&} as a ``masking'' operation: given
  values \texttt{v} and \texttt{mask}, evaluating \texttt{v \& mask}
  will only ``let through'' the bits of \texttt{v} which correspond to
  $1$ bits in \texttt{mask}; all other bits will be ``turned off''.
  For example, if you want to extract only the last three bits of a
  value \texttt{v}, you can compute \texttt{v \& 7} (since bitwise AND
  with $7 = 111_2$ will turn off all bits except the last three).

  If values are thought of as representing subsets, then \texttt{\&}
  corresponds to set intersection.
\item \texttt{|} represents bitwise logical OR.  This can be used to
  ``turn on'' certain bits: \texttt{v \& on} will result in a value
  which is the same as \texttt{v} except that the bits which are set
  to $1$ in \texttt{on} will be turned on.

  If values are thought of as representing subsets, then \texttt{|}
  corresponds to set union.
\item \verb|^| represents bitwise logical XOR.  This can be used to
  ``toggle'' bits: \verb|v ^ toggle| will result in a value which is
  the same as \texttt{v} except that the bits in positions
  corresponding to the $1$ bits in \texttt{toggle} have been flipped.

  If values are thought of as representing subsets, then \verb|^|
  corresponds to symmetric difference: \verb|a ^ b| represents the set
  of elements which are in \verb|a| or \verb|b| but not both.
\item \verb|n >> k| shifts $n$ right by $k$ bits, chopping off the
  rightmost $k$ bits.  This corresponds to (integer) division by
  $2^k$.  \verb|n << k| shifts $n$ left by $k$ bits, adding $k$ zeros
  on the right; this corresponds to multiplying by $2^k$.

  Note that right shifting uses something called \emph{sign extension}
  so that it fills in bits on the left according to whatever the
  leftmost bit was initially: a value starting with a zero bit
  (\emph{i.e.} a positive value) will have zeros filled in on the
  left, but a (negative) value beginning with a one bit will have ones
  filled in on the left.  If you don't want this (it rarely matters!)
  you can use \verb|n >>> k| which does a right shift by $k$ bits
  \emph{without} sign extension, that is, it always fills in zero bits
  on the left regardless of the initial bit of $n$.
\end{itemize}

Out of these basic operators XXX

\begin{itemize}
\item To iterate through all possible subsets of an $n$-element set
  (represented by an $n$-bit mask), just use a counter:
\begin{minted}{java}
  for (int S = 0; S < (1 << n); S++) {
      // process subset S
  }
\end{minted}
As the value of $S$ progresses from $0$ through $2^n - 1$, it will
take on every possible pattern of $n$ bits.

\item To check whether a particular bit is turned on, mask out
  everything except that particular bit and check whether the result
  is $0$.  For example, to check whether bit $j$ is set to $1$ in $S$:
\begin{minted}{java}
  if ((S & (1 << j)) != 0) ...
\end{minted}
  Be careful: the precedence of \verb|&| is actually lower than that of
  \verb|!=|, so you need a bunch of parentheses.

\item The \emph{least significant bit} (LSB) of a value \texttt{S} can
  be extracted using the expression \verb|S & (-S)|.  The result is a
  value with only a single bit set, corresponding to the LSB of $S$.
  For example, if $S = 10001011000$, then \verb|S & (-S)| will be
  $00000001000$.  It's worth taking a minute to convince yourself that
  this works, keeping in mind that to negate a value in $2$'s
  complement representation, you invert all the bits and then add one.

  One way this can be used is when iterating over all subsets of a
  set: for each subset, instead of iterating over all elements and
  checking whether each one is in the subset, one can quickly iterate
  through only the elements which are actually in the subset.  In some
  cases this can yield a big constant-factor speedup.
  \begin{minted}{java}
    for (int S = 0; S < (1 << n); S++) {
        int T = S;  // Save a copy of current subset S into T

        // Iterate through all elements of T
        while (T != 0) {
            int X = T & (-T);  // Find last element X of T
            ...                // process X
            T ^= X;            // Remove X to move on to the next element
        }
    }
  \end{minted}

  Another place this technique is used is in the implementation of
  Fenwick trees \link{sec:fenwick}.

\item The \emph{least significant zero} (LSZ) can be computed by first
  inverting all the bits, then finding the LSB.

\item There is no quick way to compute the \emph{most significant bit}
  (MSB), which amounts to finding the logarithm base 2 (rounded down).
  The simplest is to keep shifting right until reaching 1, keeping a
  count of the number of shifts required.

\item The \emph{popcount} operation counts the number of 1 bits in a
  number, and sometimes comes in useful (\inlinekattis{bits,
    iboard, enviousexponents, pebblesolitaire}).  It can be accessed
  via the \texttt{Long.bitCount(...)} or
  \texttt{Integer.bitCount(...)} functions.  Note that processors
  typically have a special popcount instruction, so this should be
  very fast---certainly much faster than manually looping through the
  bits of a number and counting how many are set to 1.

\item Iterating through all the sub-subsets of a subset
  \url{https://cp-algorithms.com/algebra/all-submasks.html}
\end{itemize}

\todo{iterating through sub-subsets}
\url{https://cp-algorithms.com/algebra/all-submasks.html}

\todo{BitSet instead of array of booleans.}

\kattis{ith}

\chapter{Geometry} \label{chap:geometry}

\kattis{alldifferentdirections, convexpolygonarea, cookiecutter,
  countingtriangles, cranes, glyphrecognition, hittingtargets,
  hurricanedanger, jabuke, janitortroubles, polygonarea, rafting}

\todo{Keep building above list---grep for geom.  Next to look at is robotprotection.}

See also list of formulas.

\todo{Points, vectors, angles.  Degrees/radians. \texttt{atan2}. Dot
  product. Rotation. Vector magnitude, norm (squared), normalize.
  Perpendicular (generate, test).}  \todo{Cross product in 2D. Signed
  area (parallelogram, triangle, Heron's formula), polygon area,
  right/left turn, inside/outside testing.}  \todo{Lines/rays (point +
  vector).  Line intersection. Segment intersection. Closest point on
  a line/segment. Point/line distance.} \todo{law of cosines.}
% -- Distance from point p to line defined by p1, p2 is absolute value
% -- of parallelogram area defined by p, p1, p2 divided by distance
% -- between p1, p2.  Proof: rearrange triangle area formula A = bh/2 to
% -- solve for h.
\todo{Convex hull.}

\url{https://codeforces.com/blog/entry/48868}

\chapter{Miscellaneous}

\section{2D grids}

2D grids/arrays (of characters, numbers, booleans\dots) are a popular
feature of many competitive programming problems.

\begin{itemize}
\item There is a trick for reading in a grid of characters which can
  save a bit of coding effort.  The ``traditional'' way to read a grid of
  characters would be something like:
  \begin{minted}{java}
    char[][] grid = new char[R][C];
    for (int r = 0; r < R; r++) {
        String line = in.nextLine();
        for (int c = 0; c < C; c++) {
            grid[r][c] = line.charAt(c);
        }
    }
  \end{minted}
  However, it is possible to assign each row of the 2D array all at
  once, like so:
  \begin{minted}{java}
    char[][] grid = new char[R][C];
    for (int r = 0; r < R; r++)
        grid[r] = in.nextLine().toCharArray();
  \end{minted}
\item In many cases the grid should be thought of as a graph where
  each cell is a vertex which is connected by edges to its neighbors.
  Note that in these cases one rarely wants to explicitly construct
  a different representation of the graph, but simply use the grid
  itself as an (implicit) graph representation.
\item It is often useful to be able to assign a unique number to each
  cell in the grid, so we can store ID numbers of cells in data
  structures rather than making some class to represent a pair of a
  row and column index.  The easiest method is to number the first row
  from $0$ to $C-1$ (where $C$ is the number of columns), then the
  second row $C$ to $2C-1$, and so on.

  \begin{tabular}{|c|c|c|c|c|}
    \hline
    $0$ & $1$ & $2$ & $\dots$ & $C-1$ \\
    \hline
    $C$ & $C+1$ & $C+2$ & $\dots$ & $2C-1$ \\
    \hline
    $2C$ & $2C+1$ & $2C+2$ & $\dots$ & $3C-1$ \\
    \hline
    $\vdots$ & $\vdots$ & $\vdots$ & $\ddots$ & $\vdots$ \\
    \hline
    $(R-1)C$ & $(R-1)C+1$ & $(R-1)C+2$ & $\dots$ & $RC-1$ \\
    \hline
  \end{tabular}

\item Using this scheme, to convert between $(r,c)$ pairs and ID
  numbers $n$, one can use the formulas
  \[ (r,c) \mapsto r \cdot C + c \qquad n \mapsto (n / C, n \% C) \]

\item To list the four neighbors of a given cell $(r,c)$ to the north,
  east, south, and west, one can of course simply list the four cases
  manually, but sometimes this is tedious and error-prone, especially
  if there is a lot of code to handle each neighbor that needs to be
  copied four times.

  Instead, one can use the following template. The idea is that
  $(dr, dc)$ specifies the \emph{offset} from the current cell $(r,c)$
  to one of its neighbors; each time through the loop we rotate it
  counterclockwise by $1/4$ turn using the mapping
  $(dr,dc) \mapsto (-dc,dr)$ (see \fulllink{chap:geometry}).

  \code{grid/Neighbors.java}
\end{itemize}

\section{Hexagonal grids} \label{sec:hex-grids}

\kattis{beehouseperimeter, honey, settlers2, beeproblem, honeyheist}

Occasionally a problem will involve a 2D grid of tiled hexagons
instead of a grid of squares.  (Typically such problems involve a
story about bees.)  They are often not too hard (\eg some kind of
straightforward application of \fulllink{sec:bfs}) other than the fact
that dealing with hexagonal grids can be annoying, unless you know a
few tricks for working with them elegantly.

\todo{Write about hexagonal grids, storage, coordinate systems, etc.}
Reference: \url{https://www.redblobgames.com/grids/hexagons/}

\section{Range queries} \label{sec:range-queries}

Suppose we have a $1$-indexed array $A[1 \dots n]$ containing some
values, and there is some operation $\oplus$ which takes two values
and combines them to produce a new value.  Given indices $i$ and $j$,
we want to quickly find the value that results from combining all the
values in the range $A[i \dots j]$, \ie $A[i] \oplus A[i+1] \oplus
\dots \oplus A[j]$.

For example, $A$ could be an array of integers, and $\oplus$ could be
$\max$, that is, we want to find the maximum value in the range
$A[i \dots j]$.  Likewise $\oplus$ could be sum, or product, or GCD.
Or $A$ could be an array of booleans, and we want to find the AND, OR,
or XOR of the range $A[i \dots j]$.

\begin{itemize}
\item For this to make sense, the combining operation must typically
  be \emph{associative}, \ie $a \oplus (b \oplus c) = (a \oplus b) \oplus
  c$.  (This is called a \emph{semigroup}.)
\item Sometimes there is also an inverse operation $\ominus$ which
  ``cancels out'' the effects of the combining operation, that is, $(a
  \oplus b) \ominus b = a$ (this is called a \emph{group}).  For
  example, subtraction cancels out addition. On the other hand, there
  is no operation that can cancel out the effect of taking a maximum.
\item If we only need to find the value of combining a \emph{single}
  range $A[i \dots j]$, then ignore everything in this section and
  simply iterate through the interval, combining all the values in
  $O(n)$ time.
\item More typically, we need to do many queries, and $O(n)$ per query
  is not fast enough.  The idea is to preprocess the array into a data
  structure which allows us to answer queries more quickly, \ie in
  $O(1)$ or $O(\lg n)$.
\item Sometimes we also need to be able to \emph{update} the array in
  between queries; in this case we need a more sophisticated query
  data structure that can be quickly updated.
\end{itemize}

Each of the below subsections outlines one approach to solving this
problem; for quick reference, each subsection title says whether an
inverse operation is required, how fast queries are, and whether the
technique can handle updates.

\subsection{Prefix scan (inverse required; $O(1)$ queries; no updates)}

In a situation where we have an inverse operation and we do not need
to update the array, there is a very simple solution.  First, make a
\emph{prefix scan array} $P[0 \dots n]$ such that $P[i]$ stores the
value that results from combining $A[1 \dots i]$.  ($P[0]$ stores the
unique ``identity'' value $a \ominus a$, \eg zero if the combining
operation is sum.)  $P$ can be computed in linear time by scanning
from left to right; each $P[i] = P[i-1] \oplus A[i]$.  Now the value
of $A[i \dots j]$ can be computed in $O(1)$ time as
$P[j] \ominus P[i-1]$. That is, $P[j]$ gives us the value of
$A[1] \oplus \dots \oplus A[j]$, and then we cancel
$P[i-1] = A[1] \oplus \dots \oplus A[i-1]$ to leave just
$A[i] \oplus \dots \oplus A[j]$ as desired.

Note that having $P[0]$ store the identity value is not strictly
necessary, but it removes the need for a special case.  If $A$ is
already $0$-indexed instead of $1$-indexed, then it's probably easier
to just put in a special case for looking up the value of $A[0 \dots
j]$ as $P[j]$, without the need for an inverse operation.

For example, suppose we are given an array of $10^5$ integers, along
with $10^5$ pairs $(i,j)$ for which we must output the sum of
$A[i \dots j]$.  Simply adding up the values in each range would be
too slow. We could solve this with the following code:

\code{range/PrefixSum.java}

More commonly, a prefix scan is a necessary first step in a more
complex solution. \kattis{divisible, dvoniz, srednji, subseqhard}

\subsection{Kadane's Algorithm}

As an aside, suppose we want to find the subsequence $A[i\dots j]$
with the \emph{biggest} sum.  A brute-force approach is $O(n^3)$:
iterate through all $(i,j)$ pairs and find the sum of each
subsequence.  Using the prefix scan approach, we can cut this down to
$O(n^2)$, since we can compute the sums of the $O(n^2)$ possible
subsequences in $O(1)$ time each.  However, there is an even better
$O(n)$ algorithm which is worth knowing, known as \emph{Kadane's
  Algorithm}.

The basic idea is simple: scan through the array, keeping a running
sum in an accumulator, and also keeping track of the biggest total
seen.  Whenever the running sum drops below zero, reset it to zero.
Below is a sample solution to \inlinekattis{commercials}.  Note that
subtracting \texttt{P} from each input is specific to the problem, but
the rest is purely Kadane's Algorithm.

\code{range/Commercials.java}

\subsection{2D prefix scan}

\todo{make pictures}

It is possible to extend the prefix scan idea to two dimensions.
Given a 2D array $A$, we create a parallel 2D array $P$ such that
$P[i][j]$ is the result of combining all the entries of $A$ in the
rectangle from the upper-left corner to $(i,j)$ inclusive.  The
simplest way to do this is to compute \[ P[i][j] = A[i][j] + P[i-1][j]
  + P[i][j-1] - P[i-1][j-1] \] Including $P[i-1][j]$ and $P[i][j-1]$
double counts all the entries in the rectangle from the upper left to
$(i-1,j-1)$ so we have to subtract them.

Given $P$, to compute the combination of the elements in some
rectangle from $(a,b)$ to $(c,d)$, we can compute \[ P[c][d] -
  P[a-1][d] - P[c][b-1] + P[a-1][b-1] \]

\inlinekattis{prozor} can be solved by brute force, but it's a nice
exercise to solve it using the above approach.

\subsection{Doubling windows (no inverse; $O(1)$ queries; no updates)}

\todo{Include link to discussion in CP3}

\subsection{Fenwick trees (inverse required; $O(\lg n)$ queries;
  $O(\lg n)$ updates)} \label{sec:fenwick}

\kattis{fenwick, supercomputer, turbo, moviecollection, dailydivision}

We can use a \emph{Fenwick tree} to query the range $A[i..j]$ (\ie get
the combination of all the values in the range $A[i] \dots A[j]$
according to the combining operation $\oplus$) in $O(\lg n)$ time.  We
can also dynamically update any entry in the array in $O(\lg n)$ time.
If dynamic updates are required and we have an invertible combining
operation, a Fenwick tree should definitely be the first choice
because the code is quite short. (Segment
trees~\link{sec:segment-trees} can also handle dynamic updates, and
work for any combining operation, even with no inverse, but the
required code is a bit longer.)

The code shown here stores \texttt{int} values and uses addition as
the combining operation, so range queries return the \emph{sum} of all
values in the range; but it can be easily modified for any other type
of values and any other invertible combining operation: change the
type of the array, change the \texttt{+} operation in the
\texttt{prefix} and \texttt{add} methods, change the subtraction
in the \texttt{range} method, and change the assignment \texttt{s = 0}
in \texttt{prefix} to the identity element instead of zero.

\begin{warning}
  Note that this \texttt{FenwickTree} code assumes the
  underlying array is $1$-indexed!
\end{warning}

\code{range/FenwickTree.java}

\begin{itemize}
\item The constructor creates a \texttt{FenwickTree} over an array of
  all zeros.
\item To create a \texttt{FenwickTree} over a given $1$-indexed array
  $A$, simply create a default tree and then loop through the array,
  calling \texttt{ft.add(i, A[i])} for each \texttt{i}.  This takes
  $O(n \lg n)$.
\item \texttt{ft.add(i, delta)} can be used to update the value at a
  particular index by adding \texttt{delta} to it.
\item If you want to simply replace the value at index $i$ instead of
  adding something to it, you could use \texttt{ft.add(i, newValue - ft.range(i,i))}.
\item \texttt{ft.range(i,j)} returns the sum $A[i] + \dots + A[j]$.
\end{itemize}

\todo{Discuss CP3 presentation of Fenwick trees; explain how Fenwick
  trees work}

\subsection{Segment trees (no inverse required; $O(\lg n)$ queries;
  $O(\lg n)$ updates)} \label{sec:segment-trees}

\todo{Segment trees.}

\section{Cycle finding} \label{sec:cycle-finding}

\todo{Floyd's algorithm, Brent's algorithm}

\chapter{Formulas} \label{chap:formulas}

\begin{itemize}
\item \textbf{Ceiling division}
  (\inlinekattis{soylent,wordcloud,amultiplicationgame}).  If
  \texttt{p} and \texttt{q} are positive values of type \texttt{int}
  or \texttt{long}, then \texttt{p/q} computes $\lfloor p/q \rfloor$,
  the quotient (rounded down).  If you want the quotient rounded
  \emph{up}, that is, $\lceil p/q \rceil$, compute
  \[ \texttt{(p + q - 1) / q}. \] Note that \texttt{-((-p)/q)} does
  not work in Java since Java truncates the result of integer division
  towards zero, instead of always taking the floor.

\item \textbf{Derangements} (\inlinekattis{secretsanta}).  The number
  of permutations of $n$ objects such that no object is left in its
  original place is
  \[ !n = n \cdot !(n-1) + (-1)^n = n! \sum_{k=0}^n \frac{(-1)^k}{k!}
    = \left[ \frac{n!}{e} \right], \]
  where $!1 = 0$, and $[x]$ denotes the closest
      integer to $x$.  The first few values of $!n$ are \[
        0,1,2,9,44,265,1854,14833,133496,1334961. \]
\item \textbf{Heron's Formula}. The area of a triangle with side
  lengths $a$, $b$, $c$ is \[ \sqrt{s(s-a)(s-b)(s-c)} \qquad
    \text{where } s = (a+b+c)/2. \]
\item \textbf{Brahmagupta's Formula} (\inlinekattis{janitortroubles}).
  The area of a quadrilateral with side lengths $a$, $b$, $c$, and
  $d$, with all vertices lying on a common circle, is
  \[ \sqrt{s(s-a)(s-b)(s-c)(s-d)} \qquad \text{where } s =
    (a+b+c+d)/2. \] This is also the maximum possible area of a
  quadrilateral with the given side lengths.
\end{itemize}

\chapter{Advanced topics} \label{chap:advanced}

This is a list of advanced topics that may eventually be included in
this document, but for now you can go read up on them if you are
interested!  (And then of course write up what you have learned for
inclusion in this document.)

\begin{itemize}
\item Chinese Remainder Theorem \kattis{heliocentric,
    generalchineseremainder}
\item Gauss-Jordan elimination (\ie row reduction \ie solving linear
  systems) (\inlinekattis{primonimo})
\item Exact Set Cover with Algorithm X/dancing links (\inlinekattis{programmingteamselection})
\item Matrix powers
  \kattis{diceandladders,driving,linearrecurrence,mortgage,overlappingmaps,squawk,timing}
\item Markov chains \kattis{lostinthewoods,gruesomecave}
\item Min cost max flow
\item Max flow with minimum and maximum capacities
\item Discrete logarithms with baby step/giant step
  (\inlinekattis{discretelogging})
\item Faster primality testing with Miller-Rabin (\eg testing with $a
  = 2, 3, 5, 7, 11, 13, 17, 19, 23, 29, 31, 37, 41$ makes it
  deterministic).
\item Divide \& conquer algorithm for counting inversions.
  \kattis{excursion, froshweek, ultraquicksort}
\item 2-SAT
\item SAT with DPLL?
\item LCA queries: Tarjan's OLCA; via RMQ; binary lifting (\inlinekattis{tourists})
\item Convolutions with Fast Fourier Transform / Number Theoretic
  Transform (\inlinekattis{tiles, aplusb, kinversions})
\end{itemize}

\appendix

\chapter{Reference} \label{chap:reference}

\section{The Alphabet}

\todo{put letter indices}

\subsection{Uppercase}

ABCDEFGHIJKLMNOPQRSTUVWXYZ

\subsection{Lowercase}

abcdefghijklmnopqrstuvwxyz

\section{ASCII}

\section{Primes}

\chapter{Resources} \label{chap:resources}

Some good resources for further learning/reference:

\begin{itemize}
\item Problems/online judges
  \begin{itemize}
  \item Of course, \href{http://open.kattis.com}{Open Kattis} has a
    collection of over 1000 great problems ranging from trivial to
    very difficult.
  \item The \href{https://uva.onlinejudge.org/}{UVa Online Judge} has
    been around much longer than Kattis and also has a huge collection
    of problems, mostly disjoint from those on Kattis.
  \item The CP3 website has a
    \href{https://cpbook.net/methodstosolve}{Methods to Solve} page
    with a huge annotated list of problems from Kattis and UVa,
    grouped by topic (corresponding to sections in CP3) with small
    hints for each one.
  \end{itemize}
\item Books
  \begin{itemize}
  \item \href{http://cpbook.net}{Competitive Programming, 3rd edition}
    (aka CP3) by Steven and Felix Halim is amazing.  Anyone serious
    about competitive programming should get a copy.
  \item
    \href{http://acm.cs.buap.mx/downloads/Programming_Challenges.pdf}{Programming
      Challenges} by Skiena and Revilla is also good.
  \end{itemize}
\item Reference
  \begin{itemize}
  \item \todo{Geeksforgeeks}
  \item \todo{Topcoder}
  \item \todo{Codeforces}
  \item \todo{cp-algorithms.com}
  \end{itemize}
\end{itemize}

\end{document}
